\documentclass[11pt]{article}
\usepackage{fontspec}
\setmainfont{Times New Roman}
\setsansfont{Arial}
\setmonofont{Consolas}
\usepackage{amsmath,amssymb,mathtools}
\usepackage{geometry}
\usepackage{microtype}
\geometry{a4paper,margin=2.28cm}
\setlength{\parindent}{1.15em}
\setlength{\parskip}{0pt}
\makeatletter
\renewcommand\footnotesize{\@setfontsize\footnotesize{7.5}{8.5}\selectfont}
\makeatother
\providecommand{\mo}{\mathfrak{o}}
\providecommand{\mM}{\mathfrak{M}}
\emergencystretch=120em
\allowdisplaybreaks
\sloppy
\hbadness=10000
\vbadness=10000
\hfuzz=2pt
\vfuzz=10pt
\raggedbottom
\begin{document}
% Унит: Папер34 сецтион15 в001. Дирецт Интерславиц Латин ауториты транслатион фром аудитед Герман соурце линес 1121--1245, сегментс P34-С0199--P34-С0211; P34-С0220/Сецтион16 ресервед фор некст транцхе. Зенодо рецорд 20822156 цхецкед фор тис унит ат 2026-06-25Т00:45З УТЦ вит но филе кеы/сизе/цхецксум цханге.
\section*{34. Хиперкомплексне величины и теорија представјень}
\emph{Продолженье и завршенье: глава III, §§15--19, и глава IV, §§20--26.}

\section*{III. поглавје. Теорија модулов и представјень}

\subsection*{§ 15. Представјеньа и модулы представјень}

Нехај $\mo$ буде колцо, а $K$ колцо с јединичным елементом. (При позднејших применьеньах $K$ јест всегда тело.)

Представјенье $n$-го степена од $\mo$ в $K$ јест хомоморфија
\[
        \mo\sim\mathfrak D,
\]
где $\mathfrak D$ јест колцо матриц $n$-го степена над $K$.

Под модулом представјеньа од $\mo$ в односу до $K$ разумеје се двојны модул $\mM$ (§1, 5), кторы јест левы $\mo$-модул и правы $K$-модул:
\[
        \mo\mM\subseteq\mM,
        \qquad
        \mM K\subseteq\mM,
\]
кторы далье јест директна сума финитно многих једнопородженых $K$-модулов:
\[
        \mM=x_1K+\cdots+x_nK;
\]
и где јединичны елемент од $K$ јест јединичны оператор.

Всакы модул представјеньа води к представјеньу. Нехај $c$ буде в $\mo$ и
\[
        cx_k=\sum_i x_i\gamma_{ik},
\]
или
\[
        c(x_1,\ldots,x_n)=(x_1,\ldots,x_n)C,
        \qquad C=(\gamma_{ik}).
\]
Тогда матрице $C$ творјат представјенье елемента $c$, бо
\[
        (b+c)x_k=\sum_i x_i(\beta_{ik}+\gamma_{ik}),
\]
\[
\begin{aligned}
        bcx_k
        &=b\sum_jx_j\gamma_{jk}
          =\sum_j b x_j\gamma_{jk}
          =\sum_j\sum_i x_i\beta_{ij}\gamma_{jk}  \\
        &=\sum_i x_i\biggl(\sum_j\beta_{ij}\gamma_{jk}\biggr),
\end{aligned}
\]
или
\[
        bc(x_1,\ldots,x_n)=(x_1,\ldots,x_n)BC.
\]

Обратно: всако представјенье $\mathfrak D$ од $\mo$ принадлежи модулу представјеньа, и то определьеној бази тог модула. Имено, под $\mM$ разумејемо целост формалных линеарных форм в $x_1,\ldots,x_n$,
\[
        y=\sum_i x_i\alpha_i.
\]
Тогда $\mM$ јест правы $K$-модул. Далье, ако елементу $c$ јест прираджена матрица $C=(\gamma_{ik})$, определимо
\[
        cx_k=\sum_i x_i\gamma_{ik},
        \qquad
        c\biggl(\sum_k x_k\alpha_k\biggr)=\sum_i x_i\sum_k\gamma_{ik}\alpha_k.
\]
Из хомоморфизмных релациј
\[
        c+d\sim C+D,
        \qquad cd\sim CD
\]
следи
\[
        (c+d)x_i=cx_i+dx_i,
        \qquad cdx_i=c(dx_i),
\]
и то само за сумы $\sum_i x_i\alpha_i$. Остале својства двојнего модула
\[
        c(y+z)=cy+cz,
        \qquad c(y\alpha)=(cy)\alpha
\]
сут тривиалне. Зато $\mM$ јест модул представјеньа, кторы по самој конструкцији точно принадлежи представјеньу $\mo\sim\mathfrak D$.

Нехај тепер $(y_1,\ldots,y_n)$ буде друга база за тој самы модул представјеньа:
\[
        (y_1,\ldots,y_n)=(x_1,\ldots,x_n)P,
        \qquad
        (x_1,\ldots,x_n)=(y_1,\ldots,y_n)P^{-1}.
\]
Ако $b\sim B$ в односу до базы $x$, тогда
\[
        b(y_1,\ldots,y_n)=b(x_1,\ldots,x_n)P
        =(x_1,\ldots,x_n)BP
        =(y_1,\ldots,y_n)P^{-1}BP.
\]
Представјеньа $b\sim B$ и $b\sim P^{-1}BP$ називајут се еквивалентне и рачуньајут се к тому самому класу представјень. Понеже всака регуларна матрица $P$ преводи базу од $\mM$ знову в базу, јест доказано:

\emph{Всакы модул представјеньа једнозначно води к определьеному класу представјень.}\textsuperscript{15а)}

{\footnotesize\noindent 15) Модулы представјень в односу до комутативных тел први раз појавјајут се у W. Крулла, \emph{Theorie und Anwendung der verallgemeinerten Abelschen Gruppen}, Sitzungsber. Heidelberger Ak. 1926, 1-а работа.\par
15а) Додано при коректури (14 априльа 1929). Како ми је сделил Б. L. van der Waerden, можно јест достати связ, независны од специалного избора базы, то јест инвариантны связ, непосреднье разделеньем појмов: линеарно трансформованье и матрица. Линеарно трансформованье јест хомоморфизм двух модулов линеарных форм; матрица јест израз, или представјенье, тог хомоморфизма при определьеном избору базы. I. Всакы модул представјеньа једнозначно прираджује области мултипликаторов $\mo$ хомоморфну систему линеарных трансформовань модула в себе. Бо по концу §2 леве мултипликаторы двојнего модула $\mM$ творјат операторне хомоморфизмы од $\mM$ в себе в односу до правых операторов (тут множенье с $K$), и то прирадженье јест хомоморфно. II. Всака система линеарных трансформовань модула линеарных форм в себе, хомоморфна до $\mo$, води к модулу представјеньа.\par}

Јасно јест: двома операторно изоморфным модулом представјень одговарја тој самы клас представјень. Але и обратно: ако два модулы представјень $(x_1,\ldots,x_n)$ и $(\bar x_1,\ldots,\bar x_n)$ творјат то само представјенье, тогда прирадженье
\[
        x_i\mapsto\bar x_i,
        \qquad
        \sum_i x_i\alpha_i\mapsto\sum_i\bar x_i\alpha_i
\]
дава изоморфизм. Бо правило множеньа јест пополно определьено матрицами $C$.

Специалны случај: ако две базы од $\mM$ творјат то само представјенье, оне преходет једна в другу чрез операторны изоморфизм од $\mM$ на себе.

Или: ако $C=PCP^{-1}$ за все $C$ и фиксовано $P$, тогда прирадженье
\[
        y_i\mapsto x_i,
        \qquad y_i=\sum_kx_k\pi_{ik},
\]
јест изоморфизм од $\mM$ на себе. Обратно: всакы изоморфизм $y_i\mapsto x_i$ двојнего модула $\mM$ на себе посредује се матрицом $P$, ктора јест заменльива со всими $C$.

Појем представјеньа можно разширити такоже на така колца, кторе сут још комутативно повезане с комутативноју областју мултипликаторов $P$ (§8). В том случају берут се само така колца-носительи представјень $K$, кторе имајут $P$ в својем центру, и од представјеньа се не требује само, же оно буде ринговы хомоморфизм $\mo\sim\mathfrak D$, але также, же оно буде операторны хомоморфизм: из $a\sim A$ имаје следовати
\[
        a\rho\sim A\rho\qquad(\rho\in P).
\]
За модулы представјень та требованьа означа додатно рачунно правило
\[
        am\cdot\rho=a(m\rho)=a\rho\cdot m.
\]
Модул и колцо тогда, како се говори, сут ``комутативно повезане с $P$''.

\end{document}

